\section{第一阶段全链路案例：从硬件事件走到内核对象}

上一节说明硬件机制如何逼出操作系统对象。本节继续用长案例把这些对象串起来。每个案例都尽量避免只给结论，而是按事件推进：硬件先发生什么，CPU 看到什么，内核保存什么状态，用户态最后观察到什么。这样写比背诵名词慢，但更接近真实调试。

\subsection{案例一：一条系统调用从用户态走到设备}

考虑用户程序执行：

\begin{lstlisting}
read(fd, buffer, 4096)
\end{lstlisting}

若 \code{fd} 指向一个冷缓存普通文件，这一行会跨过几乎所有 OS 层次。用户态先把系统调用号、fd、buffer 地址、长度放进约定寄存器，执行 syscall 指令。CPU 切到内核态，进入系统调用入口，内核保存用户寄存器形成 trap frame。

\begin{lstlisting}
user registers:
  syscall number = read
  arg0 = fd
  arg1 = user buffer pointer
  arg2 = length

CPU transition:
  switch privilege to kernel
  jump to syscall entry
  use kernel stack
  save user register frame
\end{lstlisting}

内核不能相信任何参数。fd 可能无效；buffer 可能指向未映射地址；长度可能过大；文件可能不可读；调用可能被信号打断。检查通过后，fd 表把小整数映射到 \code{struct file}，\code{struct file} 再指向 inode、文件偏移和操作函数。若 page cache 命中，内核只需从缓存页复制到用户 buffer。冷缓存时，路径继续下行。

\begin{lstlisting}
fd table -> struct file
struct file -> inode and file operations
page cache lookup(file, index)
if miss:
  filesystem maps logical file block to disk block
  block layer creates request
  driver submits DMA command to device
\end{lstlisting}

设备请求提交后，当前线程没有数据可返回。同步 \code{read} 的表现是“阻塞”，但内核内部已经变成异步：request 在设备队列里，task 在等待队列里，buffer 或 page cache folio 等待 DMA 完成。CPU 可以调度别的任务运行。

\begin{lstlisting}
current task:
  state = sleeping
  wait_reason = I/O
  wait_queue = folio waiters or request waiters

device:
  command descriptor visible through DMA
  completion will arrive by interrupt or polling
\end{lstlisting}

设备完成 DMA 后，写 completion queue 或状态寄存器，触发 MSI-X 中断。CPU 进入中断入口，驱动读取 completion，解除 DMA 映射，标记 folio uptodate，唤醒等待任务。任务重新变成 runnable，调度器稍后选中它，系统调用继续执行，把页缓存中的字节复制到用户 buffer，更新文件 offset，返回读取字节数。

\begin{longtable}{p{0.2\linewidth}p{0.34\linewidth}p{0.34\linewidth}}
\toprule
层次 & 关键对象 & 失败或延迟来源 \\
\midrule
用户态 ABI & syscall number、参数寄存器 & 参数错误、信号打断 \\
fd 层 & fd table、struct file & EBADF、权限不匹配 \\
VFS/page cache & inode、folio、xarray & cache miss、folio lock 等待 \\
文件系统 & extent、block mapping & 元数据 I/O、碎片 \\
块层 & bio、request、queue & 队列拥塞、调度延迟 \\
驱动/设备 & DMA descriptor、completion & 超时、IRQ 丢失、设备错误 \\
返回用户态 & copy\_to\_user、offset & 用户页缺页、短读 \\
\bottomrule
\end{longtable}

这个案例最重要的结论是：一次 \code{read} 不是“内核读磁盘”。它是用户态 ABI、权限、页缓存、文件系统映射、块层队列、DMA、中断、调度和用户内存复制的组合。定位慢读时也必须逐层排查：是系统调用太多，page cache miss，文件系统碎片，块队列拥塞，设备慢，还是用户 buffer 缺页。

\subsection{案例二：一次写入为什么不能保证掉电后存在}

再看：

\begin{lstlisting}
write(fd, data, len)
\end{lstlisting}

普通 buffered write 往往只把用户数据复制进 page cache，把 folio 标记为 dirty，然后返回。它没有承诺设备已经写入，更没有承诺掉电后还能读到。这个语义来自硬件速度差异和批量 I/O：如果每次小写都同步落盘，吞吐会极低，存储寿命也会受影响。

\begin{lstlisting}
buffered write:
  validate fd and user buffer
  find or allocate page cache folio
  copy bytes from user to folio
  mark folio dirty
  update inode size if needed
  return len
\end{lstlisting}

dirty folio 稍后由后台 writeback 或显式 \code{fsync} 推进到块层。若文件系统有日志，还要先记录元数据意图，再写 home location。若设备有 volatile cache，还要 flush 或使用 FUA。若目录项也改变，例如安全替换文件时使用 \code{rename}，父目录也要 \code{fsync}。

\begin{lstlisting}
durable replace:
  write temp file data
  fsync temp file
  rename temp to target
  fsync parent directory
\end{lstlisting}

这个模式看似重复，实际保护的是不同对象。第一个 \code{fsync} 保护文件数据和 inode 必要元数据；\code{rename} 改变目录项；第二个 \code{fsync} 保护目录项持久化。若省略父目录同步，崩溃后可能出现文件数据已经在盘上，但目录项没有指向它。

硬件层也有顺序问题。存储设备可能为了性能重排写入；文件系统日志必须用 barrier/flush 保证“日志提交记录”不会早于日志内容持久化。否则恢复时看见 commit，却读不到完整日志，崩溃一致性就被破坏。

\subsection{案例三：一次 page fault 如何从错误变成分配}

用户程序访问一个刚 \code{mmap} 出来的匿名页：

\begin{lstlisting}
char* p = mmap(..., PROT_READ | PROT_WRITE, MAP_PRIVATE | MAP_ANONYMOUS, ...);
p[0] = 1;
\end{lstlisting}

\code{mmap} 成功时，内核通常只建立 VMA，未必立刻分配物理页。第一次写入时，CPU 查 TLB 失败，走页表发现 PTE 不 present，于是产生 page fault。硬件把 fault address、访问类型、当前 RIP 和错误码交给内核入口。

\begin{lstlisting}
page fault handler:
  read fault address and error code
  find VMA covering address
  verify write permission
  allocate zeroed physical page
  install writable PTE
  flush or update TLB as needed
  return to faulting instruction
\end{lstlisting}

这不是“程序崩溃”。只要 VMA 合法且权限允许，内核可以修复页表，然后让同一条指令重试。重试时 TLB 能找到新 PTE，store 成功。懒分配让大虚拟地址空间变得便宜，也让实际物理内存按需使用。

COW fault 类似但更微妙。fork 后父子共享同一物理页，PTE 都是 read-only + COW。任一方写入时，page fault handler 发现这是 COW 写 fault：若页仍共享，就分配新页复制内容，再把写入方 PTE 改成 writable；若页已经独占，可以直接恢复 writable。

\begin{longtable}{p{0.2\linewidth}p{0.34\linewidth}p{0.34\linewidth}}
\toprule
fault 类型 & 内核动作 & 用户态结果 \\
\midrule
匿名懒分配 & 分配零页并安装 PTE & 指令重试成功 \\
COW 写入 & 复制共享页或恢复写权限 & 父子内容分离 \\
文件映射缺页 & 从 page cache 或存储读入 & 访问到文件内容 \\
权限错误 & 生成信号或异常 & \code{SIGSEGV} 等失败 \\
内核地址访问 & 拒绝用户访问 & 进程被终止或返回错误 \\
\bottomrule
\end{longtable}

理解 page fault 后，许多现象会变得自然：虚拟内存大于物理内存不奇怪，\code{malloc} 后首次访问慢不奇怪，fork 很快也不奇怪。真正需要警惕的是 fault 路径持锁、I/O、OOM 和 TLB shootdown 造成的长尾延迟。

\subsection{案例四：一次网络包如何变成 socket 可读}

网卡收到以太网帧时，通常通过 DMA 把数据写入预先准备好的 receive ring buffer。设备更新 completion 或 descriptor 状态，触发中断或等待 NAPI 轮询。驱动把 buffer 包装成 \code{skb}，交给网络协议栈。

\begin{lstlisting}
NIC receive:
  packet arrives on wire
  NIC chooses RX descriptor
  DMA writes packet into memory
  NIC reports completion
  driver builds skb
  protocol stack parses Ethernet/IP/TCP
  TCP queues payload to socket receive queue
  wake tasks or epoll waiters
\end{lstlisting}

这条路径有两个队列层次。硬件有 RX ring，决定设备把包放到哪里；内核有 socket receive queue，决定应用何时能读到数据。中间还有 GRO、校验和卸载、路由、Netfilter、TCP 重排、流控和拥塞控制。应用调用 \code{recv} 看到的是字节流，但内核背后处理的是包、段、序列号和状态机。

为什么高性能网络程序常用 epoll？因为 readiness 和阻塞状态要被集中管理。若 socket receive queue 为空，\code{recv} 在阻塞模式下会睡眠，在非阻塞模式下返回 \code{EAGAIN}。epoll 把“哪些 fd 可能可读”组织成事件集合，避免一个线程为大量连接逐个阻塞。

\begin{longtable}{p{0.2\linewidth}p{0.34\linewidth}p{0.34\linewidth}}
\toprule
阶段 & 关键状态 & 常见故障 \\
\midrule
RX ring & descriptor、DMA address、buffer ownership & ring 满、buffer 泄漏、DMA 错误 \\
中断/NAPI & budget、poll 状态、IRQ affinity & 中断风暴、软中断延迟 \\
协议栈 & skb、路由、TCP 状态 & 丢包、乱序、重传 \\
socket & receive queue、wait queue & backlog 满、应用读取慢 \\
用户态 & epoll interest、非阻塞循环 & 忙等、短读、背压缺失 \\
\bottomrule
\end{longtable}

这个案例说明网络不是独立于 OS 的应用层话题。它和 DMA、cache、NUMA、中断亲和性、调度延迟、内存分配都有关系。一个网络服务 p99 延迟升高，可能是 TCP 重传，也可能是 softirq 被抢占，也可能是 RX buffer 分配慢，还可能是用户态没有及时 drain socket。

\subsection{案例五：一次进程切换到底切了哪些硬件状态}

调度器选择下一个任务只是开始。真正的上下文切换要处理寄存器、栈、地址空间、FPU/SIMD 状态、调试寄存器、性能计数、CPU 本地状态和调度统计。不同架构细节不同，但核心问题相同：哪些状态是任务私有的，哪些状态是 CPU 私有的，哪些状态可以延迟保存。

\begin{lstlisting}
context switch:
  save callee-saved registers of prev
  save prev kernel stack pointer
  choose next kernel stack
  switch address space if mm differs
  handle lazy or eager FPU state
  update current task pointer
  restore next registers
  return into next task's previous kernel path
\end{lstlisting}

地址空间切换可能影响 TLB。若 x86-64 启用 PCID，旧地址空间的 TLB 项可以带标签保留；若没有标签或必须失效，切换成本更高。FPU/SIMD 状态很大，早期系统常用 lazy save，但安全问题和实现复杂性让许多系统转向 eager 策略。调度器的“公平”算法只是上层，底层切换成本会反过来影响线程数量和时间片选择。

\begin{longtable}{p{0.22\linewidth}p{0.34\linewidth}p{0.32\linewidth}}
\toprule
状态 & 为什么要处理 & 优化或风险 \\
\midrule
通用寄存器 & 任务恢复后继续计算 & 保存集合由 ABI 和入口路径决定 \\
内核栈 & 每个任务有自己的内核执行现场 & 栈溢出会破坏内核对象 \\
页表根/PCID & 不同进程地址空间不同 & PCID 降低 TLB 刷新 \\
FPU/SIMD & 用户可见寄存器很多 & lazy 策略复杂，eager 成本固定 \\
分支/cache 状态 & 微架构状态影响性能和安全 & 侧信道缓解可能增加成本 \\
\bottomrule
\end{longtable}

因此“线程很轻量”只在相对进程创建、地址空间复制而言成立。大量 runnable 线程仍会带来调度排队、cache 污染、TLB 压力和栈内存占用。OS 教材必须把调度算法和硬件切换成本一起讲，否则容易误导读者以为线程只是一个软件列表节点。

\subsection{案例六：一次设备复位为什么会影响所有等待者}

设备可能超时、固件崩溃、PCIe 链路错误或队列失去响应。驱动不能只打印错误然后继续提交请求。可靠复位要把设备状态机和内核请求状态机一起处理。

\begin{lstlisting}
device reset:
  stop accepting new requests
  mark queues frozen
  collect in-flight requests
  disable interrupts or mask vectors
  reset hardware
  rebuild rings and registers
  decide old requests: retry, fail, or complete
  unfreeze queues
  wake all affected waiters
\end{lstlisting}

这个过程解释了为什么驱动需要 request id、引用计数、超时计时器、错误码和 completion 路径。若复位时漏掉某个 in-flight 请求，等待它的进程会永久睡眠；若复位后旧 completion 迟到，可能错误完成新请求；若直接释放 DMA buffer，设备可能在复位前后继续写内存。

设备复位也是“硬件不可靠”进入 OS 设计的例子。操作系统不能假设设备总按协议完成。它必须有超时、取消、错误传播和恢复策略。高质量驱动的难点往往不在正常提交，而在这些异常路径。

\subsection{案例七：一次安全边界检查跨越了哪些硬件}

用户程序打开文件，看起来是路径权限问题，但硬件也参与了安全边界。用户态无法直接读磁盘，是因为没有设备 MMIO 权限；无法修改内核文件缓存，是因为页表隔离；无法伪造系统调用结果，是因为返回路径由内核控制；无法直接访问别的进程地址，是因为 MMU 和特权级共同限制。

\begin{lstlisting}
open(path, flags):
  syscall enters privileged mode
  kernel copies path from user memory safely
  VFS resolves path under namespace and mount rules
  DAC/capability/LSM checks decide permission
  file object is installed into fd table
  user receives only an integer handle
\end{lstlisting}

fd 是安全能力的一种形式：用户态不能直接拿 inode 指针，只能拿内核发放的整数句柄。后续 \code{read/write/ioctl} 都要通过这个句柄回到内核检查对象和操作。若 fd 泄漏到子进程，权限也随之泄漏，所以 close-on-exec 是安全边界的一部分。

\begin{longtable}{p{0.2\linewidth}p{0.34\linewidth}p{0.34\linewidth}}
\toprule
边界 & 硬件基础 & OS 策略 \\
\midrule
用户/内核 & 特权级、页表 U/S 位 & syscall、copy\_from\_user、KASLR \\
进程/进程 & 地址空间、PCID/TLB & VMA、PTE、COW、信号权限 \\
进程/文件 & 不能直接访问设备和缓存 & fd、inode 权限、LSM \\
设备/内存 & IOMMU、DMA 权限 & dma\_map/unmap、隔离域 \\
容器/宿主 & namespace 仍共享内核 & capability、seccomp、cgroup、LSM \\
\bottomrule
\end{longtable}

这说明安全不是一个章节里的附加功能，而是硬件和内核对象从一开始就共同维护的边界。后续学习 namespace、seccomp、LSM 时，应把它们看成在基础硬件边界上继续收紧权限；它们不能凭空创造隔离。

\subsection{第一阶段验收：能沿路径说清楚，而不是背定义}

完成第一阶段后，读者应该能把下面几条路径讲完整：

\begin{itemize}
\tightlist
\item 从复位向量到内核第一个 C 函数：哪些硬件已可用，哪些还不可用。
\item 从系统调用入口到返回用户态：寄存器、栈、特权级和用户指针如何处理。
\item 从虚拟地址访问到物理页：VMA、PTE、TLB、page fault 如何合作。
\item 从设备 DMA 到用户 read 返回：描述符、completion、中断、等待队列如何合作。
\item 从定时器中断到任务切换：调度账本、runqueue、上下文和地址空间如何合作。
\item 从 write 返回到 fsync 持久：page cache、writeback、日志、flush 如何合作。
\end{itemize}

如果只能背出 CPU、DMA、TLB 的定义，却不能沿这些路径说明状态变化，那还是浅层理解。第一阶段的目标就是把定义变成可追踪的状态机。
